\documentclass[10pt,a4paper]{amsart}
\usepackage{amsmath,amssymb,amsthm,graphicx}
\usepackage{verbatim}
\newenvironment{reference}{\comment}{\endcomment}
%\newenvironment{reference}{}{}
\newenvironment{slogan}{\comment}{\endcomment}
\newenvironment{history}{\comment}{\endcomment}
% For commutative diagrams we use Xy-pic
\usepackage[all]{xy}

% We use 2cell for 2-commutative diagrams.
\xyoption{2cell}
% Theorem environments.
%
\theoremstyle{plain}
\newtheorem{theorem}[subsection]{Theorem}
\newtheorem{proposition}[subsection]{Proposition}
\newtheorem{lemma}[subsection]{Lemma}

\theoremstyle{definition}
\newtheorem{definition}[subsection]{Definition}
\newtheorem{example}[subsection]{Example}
\newtheorem{exercise}[subsection]{Exercise}
\newtheorem{situation}[subsection]{Situation}

\theoremstyle{remark}
\newtheorem{remark}[subsection]{Remark}
\newtheorem{remarks}[subsection]{Remarks}

\numberwithin{equation}{subsection}

% Macros
%
\def\lim{\mathop{\mathrm{lim}}\nolimits}
\def\colim{\mathop{\mathrm{colim}}\nolimits}
\def\Spec{\mathop{\mathrm{Spec}}}
\def\Hom{\mathop{\mathrm{Hom}}\nolimits}
\def\Ext{\mathop{\mathrm{Ext}}\nolimits}
\def\SheafHom{\mathop{\mathcal{H}\!\mathit{om}}\nolimits}
\def\SheafExt{\mathop{\mathcal{E}\!\mathit{xt}}\nolimits}
\def\Sch{\mathit{Sch}}
\def\Mor{\mathop{\mathrm{Mor}}\nolimits}
\def\Ob{\mathop{\mathrm{Ob}}\nolimits}
\def\Sh{\mathop{\mathit{Sh}}\nolimits}
\def\NL{\mathop{N\!L}\nolimits}
\def\CH{\mathop{\mathrm{CH}}\nolimits}
\def\proetale{{pro\text{-}\acute{e}tale}}
\def\etale{{\acute{e}tale}}
\def\QCoh{\mathit{QCoh}}
\def\Ker{\mathop{\mathrm{Ker}}}
\def\Im{\mathop{\mathrm{Im}}}
\def\Coker{\mathop{\mathrm{Coker}}}
\def\Coim{\mathop{\mathrm{Coim}}}

% Boxtimes
%
\DeclareMathSymbol{\boxtimes}{\mathbin}{AMSa}{"02}

%
% Macros for moduli stacks/spaces
%
\def\QCohstack{\mathcal{QC}\!\mathit{oh}}
\def\Cohstack{\mathcal{C}\!\mathit{oh}}
\def\Spacesstack{\mathcal{S}\!\mathit{paces}}
\def\Quotfunctor{\mathrm{Quot}}
\def\Hilbfunctor{\mathrm{Hilb}}
\def\Curvesstack{\mathcal{C}\!\mathit{urves}}
\def\Polarizedstack{\mathcal{P}\!\mathit{olarized}}
\def\Complexesstack{\mathcal{C}\!\mathit{omplexes}}
% \Pic is the operator that assigns to X its picard group, usage \Pic(X)
% \Picardstack_{X/B} denotes the Picard stack of X over B
% \Picardfunctor_{X/B} denotes the Picard functor of X over B
\def\Pic{\mathop{\mathrm{Pic}}\nolimits}
\def\Picardstack{\mathcal{P}\!\mathit{ic}}
\def\Picardfunctor{\mathrm{Pic}}
\def\Deformationcategory{\mathcal{D}\!\mathit{ef}}
\usepackage{cleveref}
\crefname{originalreference}{original reference}{original references}
\Crefname{originalreference}{Original reference}{Original references}
\makeatletter
\newlabel{descent-lemma-locally-finite-nr-irred-local-fppf}{{descent-lemma-locally-finite-nr-irred-local-fppf}{source}{}{sourcepointer-1}{}}
\newlabel{descent-lemma-locally-finite-nr-irred-local-fppf@cref}{{[originalreference][0][]descent-lemma-locally-finite-nr-irred-local-fppf}{[1][0][]source}{}{}{}}
\makeatother
\crefname{originalreference}{original reference}{original references}
\Crefname{originalreference}{Original reference}{Original references}
\makeatletter
\newlabel{spaces-morphisms-lemma-prepare-normalization}{{spaces-morphisms-lemma-prepare-normalization}{source}{}{sourcepointer-2}{}}
\newlabel{spaces-morphisms-lemma-prepare-normalization@cref}{{[originalreference][0][]spaces-morphisms-lemma-prepare-normalization}{[1][0][]source}{}{}{}}
\makeatother
\crefname{originalreference}{original reference}{original references}
\Crefname{originalreference}{Original reference}{Original references}
\makeatletter
\newlabel{stacks-properties-lemma-quasi-compact-stack}{{stacks-properties-lemma-quasi-compact-stack}{source}{}{sourcepointer-3}{}}
\newlabel{stacks-properties-lemma-quasi-compact-stack@cref}{{[originalreference][0][]stacks-properties-lemma-quasi-compact-stack}{[1][0][]source}{}{}{}}
\makeatother
\crefname{originalreference}{original reference}{original references}
\Crefname{originalreference}{Original reference}{Original references}
\makeatletter
\newlabel{stacks-properties-lemma-type-property}{{stacks-properties-lemma-type-property}{source}{}{sourcepointer-4}{}}
\newlabel{stacks-properties-lemma-type-property@cref}{{[originalreference][0][]stacks-properties-lemma-type-property}{[1][0][]source}{}{}{}}
\makeatother
\crefname{originalreference}{original reference}{original references}
\Crefname{originalreference}{Original reference}{Original references}
\makeatletter
\newlabel{topology-lemma-surjective-continuous-irreducible-components}{{topology-lemma-surjective-continuous-irreducible-components}{source}{}{sourcepointer-5}{}}
\newlabel{topology-lemma-surjective-continuous-irreducible-components@cref}{{[originalreference][0][]topology-lemma-surjective-continuous-irreducible-components}{[1][0][]source}{}{}{}}
\makeatother
\title{Four equivalent preparations for the existence of normalizations of an algebraic stack}
\author{The Stacks Project Authors}
\date{stacks-project:stacks-morphisms.tex:local-snapshot-20261012; GNU FDL 1.2 or later}
\begin{document}
\maketitle
\noindent\textit{Editorial note.} All delivered mathematical text is the author's own, copied from the pinned Stacks Project chapter source sources/stacks-algebraic-stacks-morphisms/pinned-original-source.tex; no statement or proof step was rewritten, reordered or removed.\par\medskip
\noindent\textit{Editorial note.} Preamble provenance: the Stacks Project keeps one single shared preamble for every chapter; this card takes it directly from the staged sources/stacks-algebraic-stacks-morphisms/preamble.tex as three verbatim spans (lines 8-12, 14-18 and 154-217), so no derived file was written and nothing was added or redefined. The wrapper class and the base packages are supplied by the assembler.\par\medskip
\noindent\textit{Editorial note.} Bibliography provenance: neither the counted statement nor its proof contains a citation (this lane rescans the whole span before writing the card), so the card carries no bibliography and no entry was reconstructed.\par\medskip
\noindent\textit{Editorial note.} Selection provenance: this card was opened under the file-based control contract (CONTROL.md), with the mechanical pre-card gate run on its statement and proof interval before the card was written.\par\medskip
\noindent\textit{Allowed standard prerequisites (editorial).}\begin{itemize}
\item The corresponding conditions for morphisms of algebraic spaces. Invoked inside the counted proof (source lines 10218-10224) as cross-chapter pointers. Supplies the space-level formulation of the condition.
\item The statement that the number of irreducible components is local in the smooth topology. Invoked at the start of the counted proof (source lines 10218-10221). Makes the conditions (1) to (3) equivalent.
\item The presentation of a quasi-compact stack by an affine scheme. Invoked inside the counted proof (source lines 10228-10230) as a cross-chapter pointer. Produces the affine scheme whose finitely many components bound those of the test object.
\end{itemize}\medskip
\ifdefined\hypertarget\hypertarget{sourcepointer-1}{}\fi
\noindent\textit{Source reference.} Original reference descent-lemma-locally-finite-nr-irred-local-fppf: sources/stacks-algebraic-stacks-morphisms/pinned-original-source.tex, source lines 10219--10219. Navigation pointer only: the counted text refers to the Stacks item descent-lemma-locally-finite-nr-irred-local-fppf, which belongs to another chapter of the project and is not part of this card. The pointer carries the target's own label as its number because this chapter's staged snapshot does not contain the other chapter's tag table; its content is not reproduced here.\par\medskip
\ifdefined\hypertarget\hypertarget{sourcepointer-2}{}\fi
\noindent\textit{Source reference.} Original reference spaces-morphisms-lemma-prepare-normalization: sources/stacks-algebraic-stacks-morphisms/pinned-original-source.tex, source lines 10209--10209. Navigation pointer only: the counted text refers to the Stacks item spaces-morphisms-lemma-prepare-normalization, which belongs to another chapter of the project and is not part of this card. The pointer carries the target's own label as its number because this chapter's staged snapshot does not contain the other chapter's tag table; its content is not reproduced here.\par\medskip
\ifdefined\hypertarget\hypertarget{sourcepointer-3}{}\fi
\noindent\textit{Source reference.} Original reference stacks-properties-lemma-quasi-compact-stack: sources/stacks-algebraic-stacks-morphisms/pinned-original-source.tex, source lines 1935--1935. Navigation pointer only: the counted text refers to the Stacks item stacks-properties-lemma-quasi-compact-stack, which belongs to another chapter of the project and is not part of this card. The pointer carries the target's own label as its number because this chapter's staged snapshot does not contain the other chapter's tag table; its content is not reproduced here.\par\medskip
\ifdefined\hypertarget\hypertarget{sourcepointer-4}{}\fi
\noindent\textit{Source reference.} Original reference stacks-properties-lemma-type-property: sources/stacks-algebraic-stacks-morphisms/pinned-original-source.tex, source lines 10220--10220. Navigation pointer only: the counted text refers to the Stacks item stacks-properties-lemma-type-property, which belongs to another chapter of the project and is not part of this card. The pointer carries the target's own label as its number because this chapter's staged snapshot does not contain the other chapter's tag table; its content is not reproduced here.\par\medskip
\ifdefined\hypertarget\hypertarget{sourcepointer-5}{}\fi
\noindent\textit{Source reference.} Original reference topology-lemma-surjective-continuous-irreducible-components: sources/stacks-algebraic-stacks-morphisms/pinned-original-source.tex, source lines 10232--10232. Navigation pointer only: the counted text refers to the Stacks item topology-lemma-surjective-continuous-irreducible-components, which belongs to another chapter of the project and is not part of this card. The pointer carries the target's own label as its number because this chapter's staged snapshot does not contain the other chapter's tag table; its content is not reproduced here.\par\medskip
\begin{lemma}
\label{lemma-prepare-normalization}
Let $\mathcal{X}$ be an algebraic stack. The following are equivalent
\begin{enumerate}
\item there is a surjective smooth morphism $U \to \mathcal{X}$ where $U$
is a scheme such that every quasi-compact open of $U$ has
finitely many irreducible components,
\item for every scheme $U$ and every smooth morphism
$U \to \mathcal{X}$ every quasi-compact open of $U$ has finitely many
irreducible components,
\item for every algebraic space $Y$ and smooth morphism $Y \to \mathcal{X}$
the space $Y$ satisfies the equivalent conditions of
Morphisms of Spaces, Lemma \ref{spaces-morphisms-lemma-prepare-normalization},
and
\item for every quasi-compact algebraic stack $\mathcal{Y}$ smooth over
$\mathcal{X}$ the space $|\mathcal{Y}|$ has finitely many irreducible
components.
\end{enumerate}
\end{lemma}

\begin{proof}
The equivalence of (1), (2), and (3) follow from
Descent, Lemma \ref{descent-lemma-locally-finite-nr-irred-local-fppf},
Properties of Stacks, Lemma \ref{stacks-properties-lemma-type-property}, and
Morphisms of Spaces, Lemma \ref{spaces-morphisms-lemma-prepare-normalization}.
It is also clear from these references that condition (4)
implies condition (1). Conversely, assume the equivalent conditions
(1), (2), and (3) hold and let $\mathcal{Y} \to \mathcal{X}$ be a
smooth morphism of algebraic stacks with $\mathcal{Y}$
quasi-compact. Then we can choose an affine scheme $V$ and a surjective
smooth morphism $V \to \mathcal{Y}$ by Properties of Stacks, Lemma
\ref{stacks-properties-lemma-quasi-compact-stack}.
Since $V$ has finitely many irreducible components by (2) and since
$|V| \to |\mathcal{Y}|$ is surjective and continuous, we conclude that
$|\mathcal{Y}|$ has finitely many irreducible components by Topology, Lemma
\ref{topology-lemma-surjective-continuous-irreducible-components}.
\end{proof}

\end{document}
